% from ega3-1-fr.tex
\oldpage[III]{165}
\section*{TABLE DES MATIÈRES}

\begin{flushleft}
\textsc{Chapitre 0.} --- \textbf{Préliminaires (suite)} \dotfill 5

\medskip
\S~8. Foncteurs représentables \dotfill 5

8.1. Foncteurs représentables \dotfill 5

8.2. Structures algébriques dans les catégories \dotfill 9

\medskip
\S~9. Ensembles constructibles \dotfill 12

9.1. Ensembles constructibles \dotfill 12

9.2. Ensembles constructibles dans les espaces noethériens \dotfill 14

9.3. Fonctions constructibles \dotfill 16

\medskip
\S~10. Compléments sur les modules plats \dotfill 17

10.1. Relations entre modules plats et modules libres \dotfill 17

10.2. Critères locaux de platitude \dotfill 18

10.3. Existence d'extensions plates d'anneaux locaux \dotfill 20

\medskip
\S~11. Compléments d'algèbre homologique \dotfill 23

11.1. Rappels sur les suites spectrales \dotfill 23

11.2. La suite spectrale d'un complexe filtré \dotfill 27

11.3. Les suites spectrales d'un bicomplexe \dotfill 29

11.4. Hypercohomologie d'un foncteur par rapport à un complexe $K^\bullet$
\dotfill 32

11.5. Passage à la limite inductive dans l'hypercohomologie \dotfill 35

11.6. Hyperhomologie d'un foncteur par rapport à un complexe $K_\bullet$
\dotfill 39

11.7. Hyperhomologie d'un foncteur par rapport à un bicomplexe
$K_{\bullet\bullet}$ \dotfill 41

11.8. Compléments sur la cohomologie des complexes simpliciaux \dotfill 43

11.9. Un lemme sur les complexes de type fini \dotfill 46

11.10. Caractéristique d'Euler-Poincaré d'un complexe de modules de longueur
finie \dotfill 48

\medskip
\S~12. Compléments sur la cohomologie des faisceaux \dotfill 49

12.1. Cohomologie des faisceaux de modules sur les espaces annelés \dotfill 49

12.2. Images directes supérieures \dotfill 57

12.3. Compléments sur les foncteurs Ext de faisceaux \dotfill 60

12.4. Hypercohomologie du foncteur image directe \dotfill 62

